% Generated draft: requires mathematical and source review.
\documentclass[11pt]{article}
\usepackage[margin=25mm]{geometry}
\usepackage{fontspec}
\setmainfont{Latin Modern Roman}
\usepackage{amsmath,amssymb,mathtools,unicode-math}
\setmathfont{Latin Modern Math}
\usepackage{enumitem,needspace,xurl,hyperref}
\setlength{\parindent}{0pt}
\setlength{\parskip}{5pt}
\setlength{\emergencystretch}{3em}
\newcommand{\N}{\mathbb{N}}
\newcommand{\Z}{\mathbb{Z}}
\newcommand{\Q}{\mathbb{Q}}
\newcommand{\R}{\mathbb{R}}
\newcommand{\C}{\mathbb{C}}
\newcommand{\F}{\mathbb{F}}
\newcommand{\A}{\mathbb{A}}
\newcommand{\PP}{\mathbb{P}}
\newcommand{\E}{\mathbb{E}}
\newcommand{\eps}{\varepsilon}
\newcommand{\dd}{\,\mathrm{d}}
\begin{document}
\section*{UnsolvedMath ID: 11000087}
\textbf{Source problem code:} catalogue entry
\subsection{Definitions and mathematical statement}
Problem 4.1, ``Topological Schottky Problem,'' asks to determine the homotopy type of the source's space \(\mathcal J_g^c\), and use it to compute its ordinary and compactly supported cohomology,
\[
 H^\bullet(\mathcal J_g^c),\qquad H_c^\bullet(\mathcal J_g^c).
\]
It additionally asks to compute its intersection homology. The source gives no definitions of \(\mathcal J_g^c\), its superscript \(c\), coefficient system, or intersection-homology perversity/support convention. In particular the superscript must not be silently taken to mean a chosen closure or compactification.

The source identifies \(g=4\) as the first interesting case and states that \(\mathcal J_4^c\) is a singular divisor in \(\mathfrak h_4\). Conventionally \(\mathfrak h_g\) is Siegel upper half-space,
\[
 \mathfrak h_g=\{Z\in M_g(\C):Z^T=Z,\ \operatorname{Im}Z>0\},
\]
but the precise Jacobian-locus model, marking and topology require source confirmation.

The following context says that knowledge of \(\mathcal J_g\), \(\mathcal J_g^c\) or \(\mathcal Q_g^c\) should help compute the \(\sigma\)-invariant homology
\[
 H_\bullet(\mathcal T_g)^{\langle\sigma\rangle}.
\]
The spaces \(\mathcal Q_g^c\) and \(\mathcal T_g\), and the involution \(\sigma\), are not defined in the extract. It must not be assumed that \(\mathcal T_g\) is contractible Teichmüller space simply from the letter used. For a coefficient ring \(R\) in which \(2\) is invertible, the source gives the eigenspace decomposition
\[
 H_\bullet(\mathcal T_g;R)
 =H_\bullet(\mathcal T_g;R)^+\oplus H_\bullet(\mathcal T_g;R)^-,
\]
where \(\sigma\) acts by \(+1\) and \(-1\), respectively. This decomposition uses the projections \((1\pm\sigma)/2\) and is not asserted in characteristic \(2\) or over arbitrary integral coefficients. The proposed connection is a motivation, not an already specified isomorphism from Jacobian-locus cohomology to these eigenspaces.
\subsection{Short English statement}
Determine the topology, ordinary and compactly supported cohomology, and intersection homology of the specified Jacobian-type locus, especially its singular genus-four case. The source also seeks consequences for invariant homology of another space under an involution. The relevant spaces and comparison maps are not defined in this extract.
\subsection{Sources}
\begin{itemize}
\item Local UnsolvedMath registry, permanent ID \(11000087\), exact source code \texttt{catalogue entry}, Problem 4.1, ``Topological Schottky Problem,'' PDF page 75. No individual proposer or proposed year is supplied.
\item Farb, edited volume \emph{Problems on Mapping Class Groups and Related Topics}: \url{https://www.math.uchicago.edu/~farb/papers/mcgbook.pdf}. The extract includes chapter/page furniture ``Finiteness and Torelli spaces,'' page 69, and a footnote 7 marker without footnote text. The volume has 329 numbered items with repeated local numbering distinguished by PDF page.
\item Extraction pdf-font-structure, accessed 2026-07-29. Source evidence NEEDS\_REVIEW presents an open problem and requires a current primary-literature audit. Registry status open and classification OPEN-TRIAGE; the supplied summary explicitly says it cannot verify status. No separate authoritative definitions or computation references are supplied.
\item Dataset attribution: Ulam AI / UnsolvedMath Contributors, CC BY 4.0. Source rights NEEDS\_REVIEW: publicly accessible author/publisher copy, redistribution terms requiring review.
\end{itemize}
\end{document}
